%--------------------------------------------------------------------------------
\pagebreak
\section*{BV - Branch Vectored}
\addcontentsline{toc}{section}{BV — Branch Vectored}

%--------------------------------------------------------------------------------
\begin{iPageDescEntry}{Syntax}
\texttt{BV [.W/D/Q] [ RegX ] (  RegB ) [ , RegR ]}
\end{iPageDescEntry}

%--------------------------------------------------------------------------------
\begin{iPageDescEntry}{Format}
The \texttt{BV} instruction uses the following instruction format: \\

\begin{flushleft}
\drawInstrFormat
    {0,1,3,5,6.5,8.5,9.5,11.5,16}
    {\OpCodeGrpBR,\OpCodeBV, RegR, 0, RegB, scale, RegX, 0}
\end{flushleft}
\end{iPageDescEntry}

%--------------------------------------------------------------------------------
\begin{iPageDescEntry}{Description}
he \texttt{BV} instruction performs an unconditional region relative branch. The target address offset is formed by taking the contents of \texttt{RegX}, shifting the value according to the selected scale factor, and adding the result to \texttt{RegB}. The target address region is  the region portion of the current instruction address.

Optionally, a return link may be written to \texttt{RegR}. When specified, \texttt{RegR} receives the address of the next sequential instruction, allowing the \texttt{BV} instruction to be used for procedure calls or control-flow transfers that require a return address.

The scale factor determines the effective alignment and range of the branch offset. If no suffix is specified, the default scaling is applied.

\begin{smallGapItemize}
    \item \texttt{.W} (scale = 0) - \texttt{RegX} is shifted left by 2 bits.
    \item \texttt{.D} (scale = 1) - \texttt{RegX} is shifted left by 3 bits.
    \item \texttt{.Q} (scale = 2) - \texttt{RegX} is shifted left by 4 bits.
\end{smallGapItemize}

This mechanism enables efficient branching to word-, doubleword-, or quadword- aligned targets while keeping the instruction encoding compact.

Because \texttt{BV} is a region base relative branch, branching to a page with a different privilege level is possible. Branching from a lower privilege level to a higher one triggers an instruction protection trap. Branching from a higher privilege to a lower privilege level updates the privilege level in the processor status register accordingly. If no change is required, the privilege level remains unchanged.  
\end{iPageDescEntry}

%--------------------------------------------------------------------------------
\begin{iPageDescOpEntry}{Operation}
\begin{lstlisting}[style=iPageOpStyle]
if ( instr.[RegR] != 0 ) {
	
    RegR.[51:32] <- PSR.[IA-REG];
    RegR.[31:0]  <- addOfs32( PSR.[IA-OFS], 4 );
}

PSR.[IA-OFS] <- addOfs32( RegB, scale(RegX));
\end{lstlisting}
\end{iPageDescOpEntry}

%--------------------------------------------------------------------------------
\begin{iPageDescEntry}{Exceptions}
Taken Branch Trap\\
Instruction TLB Miss Trap
\end{iPageDescEntry}

%--------------------------------------------------------------------------------
\begin{iPageDescEntry}{Notes}
This instruction is typically used for branch tables. A branch table entry can be 1, 2, or 4 words, depending on the selected scale scheme. The branch target is computed entirely in the address-generation path and does not require additional arithmetic instructions. The optional link register write allows the instruction to serve both as a general control-flow branch and as a lightweight call mechanism.
\end{iPageDescEntry}

